\subsection{Python}
\setauthor{Elias Brandstätter}
Das Backend des Health-Monitoring-Systems wurde mit der Programmiersprache Python implementiert. Die Entscheidung für Python war dabei aus mehreren Gründen naheliegend: Einerseits entspricht Python den internen Entwicklungsrichtlinien des Unternehmens, andererseits zeichnet sich die Sprache durch eine hohe Lesbarkeit und klare Syntax aus, die es ermöglicht, auch komplexe Logik übersichtlich zu strukturieren.

Darüber hinaus verfügt Python über ein umfangreiches Ökosystem an Bibliotheken für Webentwicklung, Datenanalyse und Cloud-Integration. Gerade im Kontext dieses Projekts ist die Verfügbarkeit offizieller Clientbibliotheken für Google BigQuery von besonderer Bedeutung, da diese die Anbindung an die Datenbankinfrastruktur erheblich vereinfachen. Nicht zuletzt profitiert das Projekt von der starken Community rund um Python, die umfangreiche Dokumentationen und etablierte Best Practices bereitstellt - beides Faktoren, die die langfristige Wartbarkeit des Systems begünstigen.

\subsubsection{Hintergrund der Sprache}
\setauthor{Elias Brandstätter}
Python ist eine interpretierte, objektorientierte Programmiersprache, die ursprünglich von Guido van Rossum entwickelt und erstmals 1991 veröffentlicht wurde. Das erklärte Ziel bei ihrer Entwicklung war, eine Sprache zu schaffen, die gleichermaßen einfach zu lesen und effizient zu verwenden ist. Erreicht wird das durch eine minimalistische Syntax, die sich stark an natürlicher Sprache orientiert und dadurch die Lesbarkeit und Nachvollziehbarkeit von Code deutlich verbessert.

Im Gegensatz zu kompilierten Sprachen wird Python interpretiert ausgeführt, was die Entwicklung beschleunigt, da Änderungen am Code unmittelbar getestet werden können. Für die Entwicklung von Web-Backends und Datenverarbeitungssystemen ist diese Eigenschaft besonders vorteilhaft. \cite{python_docs, python_overview}

\subsubsection{Einsatz im Backend}
\setauthor{Elias Brandstätter}
Im Rahmen dieses Projekts übernimmt Python mehrere zentrale Aufgaben innerhalb der Backend-Architektur: die Implementierung der REST-API, die Kommunikation mit externen Schnittstellen, die Verarbeitung und Aggregation von Monitoring-Daten, Datenbankabfragen über BigQuery sowie die Logik zur Statusanalyse und Fehlererkennung.

Besonders hervorzuheben ist dabei die Fähigkeit, verschiedene externe Systeme über API-Schnittstellen anzubinden. Im Health-Monitoring-System werden Daten aus mehreren unterschiedlichen Plattformen verarbeitet - darunter App-Store-Systeme, CRM-Systeme und Projektmanagement-Tools. Python ermöglicht es, diese heterogenen Datenquellen in einer zentralen Backend-Logik zusammenzuführen und strukturiert an das Frontend weiterzugeben.

\subsection{FastAPI}
\setauthor{Elias Brandstätter}
Für die Implementierung der Backend-Schnittstellen wurde das Python-Framework FastAPI eingesetzt. FastAPI ist ein modernes, leistungsfähiges Framework zur Entwicklung von Web-APIs, das auf dem ASGI-Standard (Asynchronous Server Gateway Interface) basiert und eine schnelle sowie strukturierte Entwicklung von REST-Schnittstellen ermöglicht. Im Health-Monitoring-System stellt FastAPI die Endpunkte bereit, über die das Frontend Statusinformationen zu den überwachten Apps abruft und externe APIs angebunden werden.

\subsubsection{Hintergrund des Frameworks}
\setauthor{Elias Brandstätter}
FastAPI wurde von Sebastián Ramírez entwickelt und 2018 erstmals veröffentlicht. Das Framework macht konsequenten Gebrauch von modernen Python-Features - insbesondere Typannotationen und asynchroner Programmierung - und hebt sich dadurch von älteren Frameworks wie Flask oder Django REST Framework ab.

Ein zentrales Merkmal von FastAPI ist die Nutzung von Python-Typdefinitionen zur automatischen Validierung von Anfragen und Antworten. Request-Parameter und Response-Objekte können dadurch klar definiert werden, was die Stabilität und Vorhersehbarkeit der API erhöht. Zusätzlich generiert FastAPI automatisch eine OpenAPI-Spezifikation samt interaktiver Dokumentation, die Entwicklern direkt im Browser zur Verfügung steht und das Testen einzelner Endpunkte erheblich erleichtert. \cite{fastapi_docs}

\subsubsection{REST-API-Architektur}
\setauthor{Elias Brandstätter}
Das Backend des Health-Monitoring-Systems folgt einer klassischen REST-Architektur, bei der das System verschiedene Endpunkte bereitstellt, über die Daten abgefragt oder Aktionen ausgelöst werden können. Dazu zählen etwa das Abrufen aller Health-Checks, das Anzeigen von Details zu einem bestimmten Kunden, das Erstellen von Tasks oder Kommentaren sowie das Abfragen von Ticket-Informationen. Diese Endpunkte werden in FastAPI über Route-Deklarationen definiert, wobei jede Route einer spezifischen HTTP-Methode - etwa GET oder POST - zugeordnet ist.

Die REST-Architektur ermöglicht eine klare Trennung zwischen Frontend und Backend, sodass beide Systeme unabhängig voneinander weiterentwickelt werden können, solange die vereinbarte Schnittstellendefinition eingehalten wird.

\subsubsection{Performance und Skalierbarkeit}
\setauthor{Elias Brandstätter}
FastAPI gehört zu den performantesten Python-Webframeworks, was sich vor allem auf die native Unterstützung asynchroner Programmierung zurückführen lässt. \cite{fastapi_performance} Im Kontext dieses Projekts ist das besonders relevant, da im Zuge einer einzelnen Anfrage häufig mehrere externe APIs parallel abgefragt werden müssen. Durch die asynchrone Verarbeitung können diese Anfragen gleichzeitig ausgeführt werden, was die Gesamtantwortzeit spürbar reduziert und das System auch bei höherer Last stabil hält.

\subsection{Uvicorn}
\setauthor{Elias Brandstätter}
Zur Ausführung der FastAPI-Anwendung wird der ASGI-Server Uvicorn eingesetzt. Uvicorn ist ein leichtgewichtiger, performanter Webserver für Python-Anwendungen und übernimmt die Aufgabe, eingehende HTTP-Requests entgegenzunehmen und an die FastAPI-Anwendung weiterzuleiten. Nach der Verarbeitung durch das Framework wird die generierte Response über Uvicorn an den Client zurückgesendet. \cite{uvicorn_docs}

\subsubsection{Der ASGI-Standard}
\setauthor{Elias Brandstätter}
Uvicorn basiert auf dem ASGI-Standard (Asynchronous Server Gateway Interface), der eine moderne Alternative zum älteren WSGI-Standard darstellt. Der wesentliche Unterschied liegt in der Unterstützung asynchroner Verarbeitung: Während bei klassischen synchronen Webservern jede eingehende Anfrage einen eigenen Thread blockiert, können ASGI-basierte Server mehrere Operationen gleichzeitig verwalten, ohne dabei auf den Abschluss einer einzelnen Anfrage warten zu müssen. Das ist besonders vorteilhaft für Anwendungen, die häufig externe APIs aufrufen oder auf Datenbankabfragen warten - beides trifft auf das Health-Monitoring-System zu. \cite{asgi_spec}

\subsubsection{Einsatz im Projekt}
\setauthor{Elias Brandstätter}
Im Health-Monitoring-System fungiert Uvicorn als Laufzeitumgebung für die FastAPI-Anwendung. Der Server startet die API, stellt sie über einen definierten HTTP-Port zur Verfügung und verwaltet eingehende Anfragen, indem er diese an FastAPI weiterleitet und die Antworten zurückgibt. Durch die hohe Performance von Uvicorn kann das Backend auch bei mehreren gleichzeitigen Anfragen effizient arbeiten, ohne dass es zu merklichen Verzögerungen kommt.

\subsection{Pydantic}
\setauthor{Elias Brandstätter}
Für die Datenvalidierung und strukturierte Modellierung von Datenstrukturen wird die Python-Bibliothek Pydantic ausgewählt. Pydantic ermöglicht die Definition von Datenmodellen, die beim Empfang oder Versand von Daten automatisch validiert und bei Bedarf konvertiert werden. Im Health-Monitoring-System kommen diese Modelle vor allem bei der Verarbeitung von API-Requests und Responses zum Einsatz, wo sie sicherstellen, dass alle übermittelten Daten dem erwarteten Format entsprechen. \cite{pydantic_docs}

\subsubsection{Automatische Datenvalidierung}
\setauthor{Elias Brandstätter}
Ein zentrales Merkmal von Pydantic ist die automatische Validierung eingehender Daten. Sobald ein Client eine Anfrage an die API sendet, prüft Pydantic, ob alle erforderlichen Felder vorhanden sind und ob diese den definierten Datentypen entsprechen. Wird etwa erwartet, dass ein bestimmtes Feld eine Zeichenkette enthält oder ein Datum in einem vorgegebenen Format übergeben wird, erkennt Pydantic Abweichungen davon sofort und gibt eine entsprechende Fehlermeldung zurück. Fehlerhafte Eingaben werden damit bereits an der Systemgrenze abgefangen, bevor sie in die eigentliche Verarbeitungslogik gelangen. \cite{pydantic_validation}

\subsubsection{Datenmodelle im System}
\setauthor{Elias Brandstätter}
Pydantic verwendet sogenannte Model-Klassen, um die Struktur von JSON-Objekten zu beschreiben, die über die API übertragen werden. Im Health-Monitoring-System werden auf diese Weise Modelle für Health-Check-Ergebnisse, Kundeninformationen, Task-Objekte und Ticket-Daten definiert. Diese einheitliche Strukturierung stellt sicher, dass Daten innerhalb des gesamten Systems konsistent repräsentiert werden und keine unerwarteten Formate in der Verarbeitungslogik auftreten.

\subsubsection{Integration mit FastAPI}
\setauthor{Elias Brandstätter}
Ein besonderer Vorteil von Pydantic im Kontext dieses Projekts ist die native Integration mit FastAPI. Da FastAPI Pydantic-Modelle direkt für die Definition von Request- und Response-Schemata verwendet, entsteht eine durchgehend typsichere API, bei der Validierung, Dokumentation und Schnittstellendefinition aus einer gemeinsamen Modellbeschreibung abgeleitet werden. Das reduziert redundanten Code und trägt zur Konsistenz der gesamten Backend-Architektur bei.